\section{Related works}
\label{sec:relatedworks}
\subsection{LLM Agents in Financial Analysis}

The application of LLMs to financial tasks has grown rapidly, spanning earnings call analysis, financial question answering, and report summarisation. FinGPT \cite{yang2023fingpt} and BloombergGPT \cite{wu2023bloomberggpt} demonstrate that domain-adapted language models outperform general-purpose ones on financial benchmarks. More relevant to our work, TradingAgent \cite{tradingagents2024} and QuantAgent frame trading as a multi-step reasoning problem, delegating sub-tasks-market perception, reflection, and decision-to coordinated LLM agents. Our framework builds directly on this paradigm but extends it with vision-language chart analysis, quantitative alpha integration, and emerging-market adaptations absent from prior systems.

\subsection{Technical Analysis and Pattern Recognition}

Automating technical analysis has been explored through convolutional neural networks for candlestick pattern recognition \cite{tsantekidis2017candlestick}, reinforcement learning for indicator-based policy learning \cite{deng2016rltrading}, and hybrid rule-based systems \cite{lo2000technical}. Recent work has begun replacing handcrafted features with multimodal chart-understanding models. ChartLlama \cite{han2023chartllama} shows that vision-language models can extract actionable insights from charts without explicit feature engineering. Our system employs this capability within two dedicated vision agents, integrating it into a broader multi-agent decision pipeline rather than treating it in isolation.

\subsection{Quantitative Alpha Factor Research}

Systematic factor construction has a long tradition in quantitative finance, from foundational momentum and value factors \cite{jegadeesh1993momentum,fama1992crosssection} to the large-scale formulaic approach of Kakushadze, who codified 101 analytic alpha expressions (WorldQuant 101). Subsequent work adapts these formulas to single-asset settings \cite{kang2021singleasset} and evaluates their decay and capacity in different market regimes \cite{harvey2016,novy2016alpha}. However, traditional factor mining is notoriously susceptible to data-snooping and selection bias, pitfalls extensively documented by Harvey, Liu, and others. While modern non-LLM multimodal pipelines attempt to mitigate this by fusing factors with alternative data via standard machine learning architectures, they often lack the explicit semantic reasoning needed to contextualize anomalies. We extend this literature by constructing an 85-candidate registry-combining the full WorldQuant 101 corpus adapted to single-stock OHLCV with five proprietary formulas-and selecting the top-five factors per symbol dynamically via walk-forward Information Coefficient and Sharpe ratio ranking.

\subsection{Sentiment Analysis in Finance and Vietnamese NLP}

News sentiment has been established as a statistically significant predictor of short-term price movements \cite{tetlock2007news}, with transformer-based models such as FinBERT \cite{araci2019finbert} setting the standard for English financial text. For Vietnamese, ViSoBERT provides a BERT-based model pre-trained on large-scale Vietnamese corpora and achieves strong results on Vietnamese sentiment classification benchmarks. Financial-domain validity does not follow automatically from general-language benchmark performance. We therefore configure a hosted ViSoBERT endpoint as the primary scorer while retaining an explicit lexicon fallback, and record the actual scoring path for every newly scored article. Target and related-company sentiment are presented in a separate LLM text context and do not enter the OHLCV factor formulas or ranking. Legacy cache entries that lack scorer metadata are treated as unknown rather than retrospectively attributed to ViSoBERT.

\subsection{Trading Systems for Emerging Markets}

Research on algorithmic trading in emerging markets is comparatively sparse. Studies on the Vietnamese market have primarily examined market efficiency \cite{vo2019efficiency}, herding behaviour \cite{pham2021herding}, and the predictive power of classical technical indicators under the price-limit regime \cite{truong2020technical}, but none employs LLM-based agents. Our work develops a multi-agent LLM forecasting framework that records Vietnamese cash-equity execution and settlement assumptions separately from the domestic news ecosystem and single-stock OHLCV data characteristics; we do not claim that a universal settlement rule fixes a model horizon.

\subsection{Evaluation Standards and Agentic Benchmarks}

Recent literature highlights the critical need for rigorous, standardized evaluation in LLM-driven trading. For instance, FinVision \cite{fatemi2024finvision} shares our modular multi-agent approach and vision-based technical analysis but places much stronger emphasis on extensive ablations (e.g., of reflection modules) and comprehensive risk reporting (Sharpe, drawdowns). Furthermore, an audit-oriented survey on financial LLMs \cite{xia2026agentictrading} identifies widespread methodological gaps in time-consistent splits, transaction-cost modeling, and reproducibility. While our framework addresses costs and lookahead leakage, we recognize that our reporting detail currently falls short of the rigorous bar for full reproducibility. Finally, live/prediction-market benchmarks \cite{zhang2026predictionarena, arora2026predictionmarketbench} and forward-looking evaluation paradigms like FutureX \cite{zeng2025futurexadvancedlivebenchmark} underscore that static walk-forward backtests--while useful--must ultimately be complemented by live trials and execution realism to substantiate strong claims about agentic efficacy.
